%% Editing-only manuscript outline & status board (KEEP TO <=2 PAGES).
%% Input ONLY via \paperoutline (editing.tex); no-op in submission.tex/twocolumn.tex.
%% Status markers (shared with paper-edit skill + CLAUDE.md):
%%   \status{final} green check  = publication-ready
%%   \status{tent}  amber square = author-reviewed, keep stable (discuss before agent edits)
%%   \status{todo}  red X        = not done (agents may edit freely)
%% Six Results sections <-> six main figures. Under each: its figure, then a paragraph plan
%% (P1 = high-level overview first; P2-P4 = a few words each on where the paragraph goes).
\begingroup
\footnotesize
% Suppress the default itemize bullets -- the \status{...} symbol is the marker instead.
\renewcommand{\labelitemi}{}\renewcommand{\labelitemii}{}\renewcommand{\labelitemiii}{}
\setlength{\leftmargini}{1.1em}\setlength{\leftmarginii}{1.1em}\setlength{\leftmarginiii}{1.1em}
\setlength{\parskip}{1pt}
{\normalsize\bfseries Manuscript outline \& status}\quad
% Legend wording kept identical to notes-tex/common/tcdoc.sty so the same three
% chips read the same way in the manuscript and in the notes-tex documents.
{\scriptsize(editing view only)\quad \status{final}\,ready \enspace \status{tent}\,author-reviewed, keep stable \enspace \status{todo}\,not done, agents may edit freely}\par\vspace{1pt}
{\scriptsize\textbf{Convention:} each Results section = 3--4 paragraphs, high-level paragraph (P1) first; one main figure per section.}\par\vspace{1pt}

\textbf{Main text}
\begin{itemize}\setlength{\itemsep}{0.3pt}\setlength{\topsep}{0.5pt}
  \item \status{tent}\ \textbf{Introduction}
  \item \status{todo}\ \textbf{Results} \emph{(six sections, one per figure)}
  \begin{itemize}\setlength{\itemsep}{0.8pt}\setlength{\topsep}{0.5pt}

    \item \status{todo}\ \textbf{R1 -- Shared ontology and knowledge graph unify data}
      \begin{itemize}\setlength{\itemsep}{0pt}
        \item \status{tent}\ Fig~1 -- TorchCell + CGT overview (data a--c; model d--h)
        \item \status{todo}\ P1 \emph{(high level):} data fragmentation $\to$ one shared schema
        \item \status{todo}\ P2: Biolink ontology, ingest-time validation, Neo4j KG, datasets-as-queries
        \item \status{todo}\ P3: reference-cell $+$ perturbation factorization; join/aggregation shrinks labels
        \item \status{todo}\ P4: information budget ($10^{10}$ vs $10^{8}$ bits); cut at $H$; baselines saturate fitness
      \end{itemize}

    \item \status{todo}\ \textbf{R2 -- CGT predicts trigenic gene--gene interactions (state of the art)}
      \begin{itemize}\setlength{\itemsep}{0pt}
        \item \status{todo}\ Fig~2 -- gene--gene interactions + iBioFoundry design loop
        \item \status{todo}\ P1 \emph{(high level):} interactions are the hard structured label; CGT in one line
        \item \status{todo}\ P2: trigenic interaction defined (peel off lower-order effects)
        \item \status{todo}\ P3: results -- SOTA vs Yeast9 FBA / DCell / DANGO ($r{=}0.454$)
        \item \status{todo}\ P4: data/model scaling, graph-reg ablation, accuracy vs $|\tau|$
      \end{itemize}

    \item \status{todo}\ \textbf{R3 -- One embedding, many phenotypes}
      \begin{itemize}\setlength{\itemsep}{0pt}
        \item \status{todo}\ Fig~3 -- expression, morphology, fitness
        \item \status{todo}\ P1 \emph{(high level):} one latent embedding generalizes across phenotypes
        \item \status{todo}\ P2: expression prediction ($r{\approx}0.54$) + diagnostics
        \item \status{todo}\ P3: single-knockout morphology ($r{\approx}0.62$)
        \item \status{todo}\ P4: fitness + multitask sharing; why one embedding helps
      \end{itemize}

    \item \status{todo}\ \textbf{R4 -- Natural genetic variation vs.\ model-designed perturbations}
      \begin{itemize}\setlength{\itemsep}{0pt}
        \item \status{todo}\ Fig~4 -- natural genetic variation
        \item \status{todo}\ P1 \emph{(high level):} natural variation vs designed perturbations
        \item \status{todo}\ P2: natural strain variation across genome + environment
        \item \status{todo}\ P3: model-designed (in-silico) gene perturbations
        \item \status{todo}\ P4: comparison; what design adds beyond the natural range
      \end{itemize}

    \item \status{todo}\ \textbf{R5 -- Drug exposure and representation robustness}
      \begin{itemize}\setlength{\itemsep}{0pt}
        \item \status{todo}\ Fig~5 -- drug exposure + robustness \emph{(paper tag \texttt{r-robust})}
        \item \status{todo}\ P1 \emph{(high level):} operator extends to drug/env; representation is robust
        \item \status{todo}\ P2: drug-exposure prediction
        \item \status{todo}\ P3: robustness -- stable across splits; embedding + graph ablations
        \item \status{todo}\ P4: graceful degradation as supervision shrinks
      \end{itemize}

    \item \status{todo}\ \textbf{R6 -- Metabolism and strain design}
      \begin{itemize}\setlength{\itemsep}{0pt}
        \item \status{todo}\ Fig~6 -- metabolism + metabolic engineering \emph{(paper tag \texttt{r5})}
        \item \status{todo}\ P1 \emph{(high level):} ground the representation in metabolism; recommend targets
        \item \status{todo}\ P2: metabolism integration (Yeast9 / gene--reaction)
        \item \status{todo}\ P3: strain-design recommendations + iBioFoundry DBTL loop
        \item \status{todo}\ P4: inference at scale
      \end{itemize}

  \end{itemize}
  \item \status{todo}\ \textbf{Discussion}
  \item \status{todo}\ \textbf{Methods} \emph{(subsection status inline in \texttt{methods.tex})}
\end{itemize}

\textbf{Supplementary Information}
\begin{itemize}\setlength{\itemsep}{1pt}\setlength{\topsep}{0.5pt}

%% Numbers below are LIVE cross-references (\suppfig/\supptab/\ref), never hand-typed: the board
%% cannot drift out of sync with the SI again. Nature style -- "Supplementary Fig. 1", no "S" prefix.
  \item \status{tent}\ \textbf{Note~\ref{note:equivalence} -- Functional equivalence of the perturbation operator}
    \begin{itemize}\setlength{\itemsep}{0pt}
      \item \status{todo}\ \suppfig{fig:pert-equivalence-empirical} -- empirical equivalence (CGT- vs GNN-style, KO fitness)
      \item \status{tent}\ Setting · Concrete example · Origin of approximation · Precision · Relation to prior work · Empirical validation
    \end{itemize}
  \item \status{tent}\ \textbf{Note~\ref{note:graph-attention} -- Graph-regularized attention contains graph attention}
    \begin{itemize}\setlength{\itemsep}{0pt}
      \item \status{todo}\ \suppfig{fig:graphreg-empirical} -- graph-reg relevance ($\lambda$ sweep)
      \item \status{tent}\ Setting · Precision · Empirical validation
    \end{itemize}
  \item \status{tent}\ \textbf{Note~\ref{note:static-graph} -- Static-graph assumption and learning the perturbation}
    \begin{itemize}\setlength{\itemsep}{0pt}
      \item \status{tent}\ \supptab{tab:compute} -- per-batch compute (two routes)
      \item \status{tent}\ Static-graph assumption · Limitations · Single general data object · Compute \& complexity · Empirical cost · Trade-offs
    \end{itemize}
  \item \status{tent}\ \textbf{Note~\ref{note:gi-from-fitness} -- Gene interactions as a learned function of fitness}
    \begin{itemize}\setlength{\itemsep}{0pt}
      \item \status{todo}\ \suppfig{fig:fit-gi-emergent} -- do interactions emerge from fitness?
      \item \status{tent}\ Setting · Toward a formal statement · What would be surprising · Implication for metabolism · The atomic unit of data
    \end{itemize}
  \item \status{todo}\ \textbf{Note~\ref{note:info-accounting} -- Persistent entities and contingent observations} \emph{(theory of the cut; Fig.~1c)}
    \begin{itemize}\setlength{\itemsep}{0pt}
      \item \status{todo}\ \supptab{tab:entity-corpora} -- persistent-entity corpora, measured from the public archives by \texttt{experiments/016-information-accounting}
      \item \status{todo}\ Setting · The measure (Prop.: compressed length bounds, not entropy) · The two corpora · Two asymmetries · Consequence · Objection · Precision
      \item \status{todo}\ \emph{Open:} contingent-corpus gzip figure still computed on gilahyper -- commit that script and regenerate \supptab{tab:datasets} with a compressed-size column
    \end{itemize}
  \item \status{todo}\ \textbf{Note~\ref{note:classical-ml} -- Classical-ML case study: representation \& pooling} \emph{(empirical counterpart of Note~\ref{note:info-accounting})}
    \begin{itemize}\setlength{\itemsep}{0pt}
      \item \status{todo}\ \suppfig{fig:classical-ml} -- classical-ML baselines (compose from \texttt{*\_spearman\_spearman\_*\_palette} bars + \texttt{node\_embedding\_performance\_*\_palette}) · \supptab{tab:classical-ml-results}--\ref{tab:classical-ml-pearson} (Spearman, MSE, Pearson)
      \item \status{tent}\ Setup · Representation is identity · The barrier is pooling
    \end{itemize}
  \item \status{todo}\ \textbf{Note~\ref{note:graphs} -- Gene--gene graphs used as attention priors} \emph{(the nine graphs of experiment 010; STRING releases)}
    \begin{itemize}\setlength{\itemsep}{0pt}
      \item \status{todo}\ \suppfig{fig:graph-priors} -- sizes, degree, Jaccard, containment, multiplicity, structure · \suppfig{fig:graph-priors-2} -- shared pairs, hubs, regulatory vs TFLink, STRING v9.1/v11.0/v12.0 drift (\texttt{experiments/010-kuzmin-tmi/scripts/graph\_statistics.py}) · \supptab{tab:databases}--\ref{tab:string-versions}
      \item \status{todo}\ Sources · Size, overlap, and structure · Hubs and the two transcription-factor graphs · STRING releases
    \end{itemize}
  \item \status{todo}\ \textbf{Note~\ref{note:dango-repro} -- Constructing DANGO in TorchCell and reproducing its published result} \emph{(Kuzmin 2018 screen alone; 19 runs at $r=0.42$; the lambda fall-through)}
    \begin{itemize}\setlength{\itemsep}{0pt}
      \item \status{todo}\ \suppfig{fig:dango-repro} -- model in the shared notation, decreased zeros, best $r$ by release and schedule, curves (\texttt{experiments/005-kuzmin2018-tmi/scripts/dango\_construction\_si.py}, \texttt{dango\_string\_version\_sweep.py}) · \supptab{tab:dango-string-versions}
      \item \status{todo}\ Model in the manuscript's notation · Per-channel zero weights · Dataset, result, and what is not matched \quad \status{todo}\ \emph{Open:} DANGO reference key \texttt{zhangDANGOPredictingHigherorder2020} not yet in \texttt{references.bib}
    \end{itemize}
  \item \status{todo}\ \textbf{Note~\ref{note:dango-full} -- DANGO on the full trigenic dataset} \emph{(the Fig.~2d value; canonical statement of the 006-vs-010 build mismatch)}
    \begin{itemize}\setlength{\itemsep}{0pt}
      \item \status{todo}\ \suppfig{fig:dango-full} -- curves, best $r$ by release, epochs to $r\ge0.35$; panels d--f are placeholders needing a GPU run (\texttt{experiments/010-kuzmin-tmi/scripts/dango\_full\_dataset\_si.py}) · \supptab{tab:dango-full-runs}
      \item \status{todo}\ Dataset builds · Training and result · Release and dataset effects \quad \status{todo}\ \emph{Open:} no DANGO or DCell run on the experiment-010 build (the like-for-like baseline)
    \end{itemize}
  \item \status{todo}\ \textbf{Note~\ref{note:dcell-model} -- DCell in TorchCell: the ontology-structured model} \emph{(GO DAG as the model graph; the rebuilt 2,655-subsystem ontology)}
    \begin{itemize}\setlength{\itemsep}{0pt}
      \item \status{todo}\ \suppfig{fig:dcell-model} -- schematic + subsystems per stratum, genes per subsystem, subsystems per gene (\texttt{experiments/006-kuzmin-tmi/scripts/dcell\_model\_go\_stats.py}) · \supptab{tab:dcell-go-filter}
      \item \status{todo}\ The graph is the ontology · Building and measuring the ontology \quad \status{todo}\ \emph{Open:} DCell (\texttt{maUsingDeepLearning2018}) and Gene Ontology reference keys not yet in \texttt{references.bib}
    \end{itemize}
  \item \status{todo}\ \textbf{Note~\ref{note:dcell-training} -- DCell training cost, instability, and the checkpoint-based baseline} \emph{(one 30.8-day run; the Fig.~2d error bar is checkpoint spread, not replicates)}
    \begin{itemize}\setlength{\itemsep}{0pt}
      \item \status{todo}\ \suppfig{fig:dcell-training} -- validation curve with the three checkpoints, losses, GPU-hours and throughput vs DANGO and CGT, speed-up stages; panel e is a placeholder needing a profiler run (\texttt{experiments/006-kuzmin-tmi/scripts/dcell\_training\_wandb.py}) · \supptab{tab:dcell-training-checkpoints}--\ref{tab:dcell-training-cost}
      \item \status{todo}\ Training run and checkpoint statistic · Cost and where the time goes
    \end{itemize}

  \item \status{tent}\ \textbf{General SI figures} (follow the Notes)
    \begin{itemize}\setlength{\itemsep}{0pt}
      \item \status{tent}\ \suppfig{fig:overview} -- TorchCell overview schematic \quad \status{tent}\ \suppfig{fig:data-model} -- ontology data model (hourglass)
      \item \status{tent}\ \suppfig{fig:conversion-dedup-agg} -- KG normalization (convert/dedup/aggregate) \quad \status{tent}\ \suppfig{fig:dna-selection} -- DNA sequence selection
      \item \status{tent}\ \suppfig{fig:dna-tokenization} -- DNA tokenization \quad \status{todo}\ \suppfig{fig:expr-diag} -- expression \& multitask diagnostics
      \item \status{todo}\ \suppfig{fig:inference-scale} -- inference at scale \quad \status{tent}\ \suppfig{fig:profile-enrichment} -- guilt-by-association annotation
    \end{itemize}
  \item \status{tent}\ \textbf{SI tables}: \status{tent}\,\supptab{tab:compute} compute \quad \status{todo}\,\supptab{tab:entity-corpora} entity corpora \quad \status{todo}\,\supptab{tab:classical-ml-results}--\ref{tab:classical-ml-pearson} classical-ML \quad \status{tent}\,\supptab{tab:datasets} datasets \quad \status{todo}\,\supptab{tab:databases} graphs (now script-generated; counts changed) \quad \status{todo}\,\supptab{tab:string-versions} STRING releases \quad \status{todo}\,\supptab{tab:dango-string-versions} DANGO by release \quad \status{todo}\,\supptab{tab:dango-full-runs} DANGO full-dataset runs \quad \status{todo}\,\supptab{tab:dcell-go-filter} DCell GO filter \quad \status{todo}\,\supptab{tab:dcell-training-checkpoints} DCell checkpoints \quad \status{todo}\,\supptab{tab:dcell-training-cost} training cost
\end{itemize}
\endgroup
